\chapter{Discussion}
\label{sec:discussion}

\section{Discussion of Evaluation Results}
\label{sec:results-discussion}


\section{Rust versus Haskell}
\label{sec:rust-vs-haskell}

\subsection{Language Characteristics of the Reference Implementation}
\label{sec:haskell-language-characteristics}

The reference material this thesis translates from -- both
flops14-extended's own presentation \citep{sulzmannlu2014} and
Sulzmann's fuller teaching-page Haskell transcript \citep{sulzmann_slides}
-- leans on several Haskell features with no direct Rust counterpart,
each already surfaced individually at the point in this thesis where it
first mattered; this section collects them.

\paragraph{Algebraic data types with blanket \texttt{deriving}.} The
teaching page's own regex type,
\begin{verbatim}
data R = Phi | Eps | L Char | Seq R R | Choice R R | Star R
  deriving (Show, Eq, Ord)
\end{verbatim}
derives structural equality, ordering, \emph{and} printing in one clause,
whether or not a given piece of code actually uses all three. Section~\ref{sec:regex-rust-implementation}'s
\texttt{Regex} derives only \texttt{Debug}, \texttt{Clone},
\texttt{PartialEq}, \texttt{Eq}, and \texttt{Hash} -- no \texttt{Ord} --
because nothing in this thesis's Rust port ever compares two
\texttt{Regex} values for ordering; Rust's per-trait, opt-in
\texttt{derive} makes that omission a visible, deliberate fact about the
codebase, where the reference's blanket \texttt{deriving} leaves it
implicit.

\paragraph{GADT syntax used without a type index.} The reference's parse
tree type is written as a GADT,
\begin{verbatim}
data U where
  Nil :: U
  Empty :: U
  Letter :: Char -> U
  LeftU :: U -> U
  RightU :: U -> U
  Pair :: (U,U) -> U
  List :: [U] -> U
\end{verbatim}
but no constructor's return type is ever refined to anything other than
plain \texttt{U} -- this is an ordinary sum type written in GADT syntax,
not a use of GADTs' actual type-level power. Section~\ref{sec:parsetree-and-flatten}'s
\texttt{ParseTree} enum needs nothing beyond an ordinary Rust
\texttt{enum} to match it.


\subsection{What the Reference Provides vs. What It Leaves Open}
\label{sec:reference-scope}

\subsection{Translation Strategy}
\label{sec:translation-strategy}

\subsection{Translation Insights}
\label{sec:translation-insights}

\section{Mismatch Analysis via Derivative Tracing}
\label{sec:mismatch-analysis}

\subsection{The Gap}
\label{sec:diagnostics-gap}


\subsection{Deriving Diagnostics for Free from the Algorithm Itself}
\label{sec:diagnostics-for-free}

\subsection{Design}
\label{sec:diagnostics-design}

\texttt{DiagLevel} (\texttt{src/diagnostics/mod.rs}) has four values --
\texttt{Off} (0), \texttt{Basic} (1, regex/input/match/tree/error caret),
\texttt{Verbose} (2, \texttt{Basic} plus timing, expression count, and
construction steps), and \texttt{Debug} (3, the full derivation trace,
written to a file if \texttt{REGEX\_DIAG\_REPORT} is set) -- selected via
the CLI's \texttt{--diag} flag or the \texttt{REGEX\_DIAG} environment
variable, and applies only to \texttt{parse}, not \texttt{match}: a
Boolean matcher has no parse tree, and hence nothing at levels 1--3 to
report beyond what the trace itself would show, which level 3 already
provides directly. Each of the four parser variants renders its own
trace independently rather than sharing one generic renderer --
\texttt{diag\_2/\{deriv\_std,deriv\_bc,pderiv\_std,pderiv\_bc\}.rs} and
the analogous four files under \texttt{diag\_3/} -- since what is worth
showing at each level differs by construction (a bit-coded parser's trace
is bit-strings and \texttt{ARegex} states; the tree-injection parsers'
traces are \texttt{Regex} derivative sequences and, at level 3, the
backward injection steps that reconstruct a tree from them).
\texttt{docs/DIAG.md} documents concrete output at every level, for every
parser, on both successful and failing matches.

\subsection{Scope and Limits}
\label{sec:diagnostics-scope-limits}

\section{Optimization Challenges and Insights}
\label{sec:optimization-challenges}


\section{Future Work}
\label{sec:future-work}

\subsection{Paper-Suggested Extensions}
\label{sec:paper-suggested-extensions}

\subsubsection{Tagged NFA for Submatching}
\label{sec:tagged-nfa}

\subsubsection{Space-Efficient Parsing}
\label{sec:space-efficient-parsing}

\subsubsection{DFA/NFA Hybrid Abstraction}
\label{sec:dfa-nfa-hybrid}

\subsection{Further Performance Optimizations}
\label{sec:further-performance-optimizations}

\subsubsection{Character Classes in Place of Literal Alternation}
\label{sec:char-class-optimization}

\subsubsection{Structural Sharing via \texttt{Rc}}
\label{sec:rc-sharing-optimization}

\subsubsection{Memoizing \texttt{deriv}/\texttt{nullable}}
\label{sec:memoization-optimization}

\subsubsection{Benchmarking the Matcher, Not the Parser, for Membership}
\label{sec:matcher-benchmark-fairness}

\subsection{A Full Mismatch Diagnostics System}
\label{sec:mismatch-diagnostics-future}
